\subsection{Molecular simulation of heat treatment}
\label{sec:md_results}
The results of the MD simulation are summarized in \cref{fig:diffusion}. The simulated heat treatment was conducted at \SI{895.15}{\kelvin}, a temperature just below Al's melting point of \SI{933.5}{\kelvin} and significantly lower than Ti's melting point of \SI{1941}{\kelvin}. It should be noted that under these conditions, Al atoms are likely to exhibit more dynamic activity than Ti atoms. For Ti atoms in the pure Ti phase to disengage from their local environment and diffuse, a fundamental modification and activation of the Ti surface is necessary.

\cref{fig:atomic_diffusion} depicts the dynamic atom movements during the heat treatment process. It highlights a significant diffusion of Al atoms from the pure Al phase into the TiAl\textsubscript{3} grain boundaries ($0 < z < 40$ \AA), likely a result of pre-melting as evidenced by the observation of liquid Al at the interfaces. An apparent accumulation of Al atoms near the pure Ti phase is also observed, as indicated by an increased density near $z = 40$ \AA. However, the figure illustrates limited diffusion of Ti atoms, whether they initially reside in the pure Ti or the TiAl\textsubscript{3} phase. The Ti surface remains largely inactive in the simulation and doesn't exhibit substantial diffusion even when Al atoms infiltrate the Ti/TiAl\textsubscript{3} interface. Despite the co-existence of Al diffusion into the Ti side and Ti diffusion into the Al side, the former dominates. Ti atoms from the pure Ti phase rarely migrate into the Al side or into the TiAl\textsubscript{3} grain boundaries, suggesting that Ti atom diffusion might only be noticeable at very extended time-scales due to long activation times. Moreover, Al atoms initially in the TiAl\textsubscript{3} phase show negligible diffusion, likely due to the stability of the TiAl\textsubscript{3} lattice structure at the heat treatment temperature.

\begin{figure}[htbp]
    \def\subplotwidth{0.85\textwidth}
    \def\verticalskip{0.5\baselineskip}
    \newcommand{\subplot}[2]{
        \sidesubfloat[]{
            \includegraphics[width=\subplotwidth]{#1}\label{#2}
        }
    }
    \subplot{figs/md_simulation_results/atomic_diffusion}{fig:atomic_diffusion}
    \\
    [\verticalskip]
    \subplot{figs/md_simulation_results/Al_diffusion}{fig:Al_diffusion}
    \caption{The dynamic evolution of atoms during heat treatment simulation. (a) shows the evolution of density profiles of all the atoms along the z-direction during heat treatment simulation. The shaded region corresponds to the TiAl\textsubscript{3} grain layer. (b) depicts the variation of the z-positions of Al atoms from the pure Al layer. The colorbar corresponds to the range of the TiAl\textsubscript{3} grain layer, with the lower limit representing the bottom and the upper limit representing the top of the layer. The entire pure Al region lies below the TiAl\textsubscript{3} layer and is shown in dark blue. The "sliced" visualizations in the bottom row present a top-view of the Al atoms with the atoms located above z-position $20$ removed, providing insight into the configuration of Al atoms in the TiAl\textsubscript{3} grain boundaries.}
    \label{fig:diffusion}
\end{figure}

Therefore, the most predominant process is the diffusion of Al atoms from the pure Al layer towards the Ti layer via the TiAl\textsubscript{3} grain boundaries, followed by the accumulation of Al atoms at the Ti/TiAl\textsubscript{3} interface. \cref{fig:Al_diffusion} specifically focuses on this phenomenon. As demonstrated in the figure, during the heat treatment, the Al atoms diffuse along the TiAl\textsubscript{3} grain boundaries, causing them to slowly broaden. Upon reaching the Ti/TiAl\textsubscript{3} interface, the Al atoms exhibit a tendency to accumulate. This result is consistent with literature that the Ti-Al alloy forming mostly happens in the vicinity of pure Ti layer.\cite{thiyaneshwaran2018nucleation, mirjalili2013kinetics} The findings indicate that in additive manufacturing, post-processing a manufactured part with remaining pure Ti phase embedded in the Al matrix for full reaction of Ti with Al could be challenging. This difficulty arises due to the considerable disparity between the diffusion rates of Al and Ti, which may lead to the formation of Kirkendall voids at the Al/TiAl\textsubscript{3} interface. The issue can potentially be mitigated by the recently established \textit{in situ} reactive printing technique in additive manufacturing.\cite{tian2022situ} This method leverages on the rapid cyclic thermal conditions of metal additive manufacturing to instantaneously form reinforcing reaction product during the printing process.

\subsection{Distribution of states}
\label{sec:distribution_of_states}
The MD simulation result shows that the diffusion of pure Al atoms is dominant in this system. Therefore, we focus our analysis on the dynamic behavior of these atoms.

The results for the Al-substrated and Ti-substrated systems at \SI{1033.5}{\kelvin} are presented in \cref{fig:Al-based} and \cref{fig:Ti-based}, respectively. In both systems, three Markov states of the dynamics of Al are observed:
\begin{itemize}
    \item state 1 (solid Al): a layer of Al atoms that are near the substrates
    \item state 2 (interfacial Al): Al atoms at the interfacial region
    \item state 3 (liquid Al): Al atoms in the mixture that is far away from the substrates
\end{itemize}
% with the corresponding atoms consisting of a layer of Al atoms near the substrates (state 1), a interfacial layer (state 2), and a region further away from the substrate where the dynamics are not affected by it (state 3). 
Remarkably, the MSM was trained solely on local information and not on the exact atomic positions, yet it accurately captures the spatial distribution of the learned states on a global scale.

\begin{figure}[htbp]
    \def\subplotheight{0.18\textwidth}
    \def\verticalskip{0.5\baselineskip}
    \newcommand{\subplot}[2]{
        \sidesubfloat[]{
            \includegraphics[height=\subplotheight]{#1}\label{#2}
        }
    }

    \subplot{figs/distribution_and_probability/Al-Al100Ti0-Al-7-distribution-A_exam.png}{fig:Al-distr_sub0}
    \subplot{figs/distribution_and_probability/Al-Al100Ti0-Al-7-probability-A_exam.png}{fig:Al_prob_sub0}
    \\
    % [\verticalskip]

    \subplot{figs/distribution_and_probability/Al-Al95Ti05-Al-9-distribution-A_exam.png}{fig:Al-distr_sub1}
    \subplot{figs/distribution_and_probability/Al-Al95Ti05-Al-9-probability-A_exam.png}{fig:Al_prob_sub1}
    \\
    % [\verticalskip]

    \subplot{figs/distribution_and_probability/Al-Al90Ti10-Al-7-distribution-A_exam.png}{fig:Al-distr_sub2}
    \subplot{figs/distribution_and_probability/Al-Al90Ti10-Al-7-probability-A_exam.png}{fig:Al_prob_sub2}
    
    \caption{The distribution of states of Al atoms in Al-substrated system at 1033.5K. The substrates are omitted. (a), (c) and (e) depict the spatial distribution of probabilities of Al atoms being in different states for composition Al$100\%$Ti$0\%$, Al$95\%$Ti$5\%$ and Al$90\%$Ti$10\%$, respectively. (b), (d) and (f) are the corresponding probability distribution as a function of the z position.}
    \label{fig:Al-based}
\end{figure}

\begin{figure}[htbp]
    \def\subplotheight{0.18\textwidth}
    \def\verticalskip{0.5\baselineskip}
    \newcommand{\subplot}[2]{
        \sidesubfloat[]{
            \includegraphics[height=\subplotheight]{#1}\label{#2}
        }
    }

    \subplot{figs/distribution_and_probability/Ti-Al100Ti0-Ti-1-distribution-A_exam.png}{fig:Ti-distr_sub0}
    \subplot{figs/distribution_and_probability/Ti-Al100Ti0-Ti-1-probability-A_exam.png}{fig:Ti_prob_sub0}
    \\
    % [\verticalskip]

    \subplot{figs/distribution_and_probability/Ti-Al95Ti05-Ti-9-distribution-A_exam.png}{fig:Ti-distr_sub1}
    \subplot{figs/distribution_and_probability/Ti-Al95Ti05-Ti-9-probability-A_exam.png}{fig:Ti_prob_sub1}
    \\
    % [\verticalskip]

    \subplot{figs/distribution_and_probability/Ti-Al90Ti10-Ti-7-distribution-A_exam.png}{fig:Ti-distr_sub2}
    \subplot{figs/distribution_and_probability/Ti-Al90Ti10-Ti-7-probability-A_exam.png}{fig:Ti_prob_sub2}
    
    \caption{The distribution of states of Al atoms in Ti-substrated system at 1033.5K. The substrates are omitted. (a), (c) and (e) depict the spatial distribution of probabilities of Al atoms being in different states for composition Al$100\%$Ti$0\%$, Al$95\%$Ti$5\%$ and Al$90\%$Ti$10\%$, respectively. (b), (d) and (f) are the corresponding probability distribution as a function of the z position.}
    \label{fig:Ti-based}
\end{figure}

From \cref{fig:Al-based} and \cref{fig:Ti-based} we can see that the MSMs are more confident in classifying Al atom states as the percentage of Ti in the mixture increases. For example, in Al-based systems, the MSM's confidence in identifying atoms in the middle of the mixture as liquid Al increases with the concentration of Ti in the mixture. Ti-based systems display a similar state distribution, except that in \cref{fig:Ti-based}, composition $\mathrm{Ti}0\%\mathrm{Al}100\%$, the MSM struggles to differentiate between state 2 and state 3 for many atoms. This could possibly be due to that the effect of the Ti substrate on the dynamics of Al atoms has a relatively shorter range compared to the Al substrate, so, when the composition of the mixture is as simple as pure Al, state 2 and 3 are strongly overlapping in the space.

State distributions at \SI{983.5}{\kelvin} are detailed in \cref{asec:distribution_at_983.5K}, showing similarity to the results at \SI{1033.5}{\kelvin}. Seemingly, the MSMs demonstrate greater confidence in classifying Al atom states at the lower temperature, possibly due to more distinct transition modes between states. As a supplementary analysis, we compared the distribution of Markov states with the potential energy distribution of Al atoms in \cref{asec:potential_energy_distribution}. It was noted that the two distributions do not always coincide, underlining that an understanding of the dynamics cannot be derived from the analysis of potential energy due to its inherent static nature.

\subsection{Transitions between states}
\label{sec:transitions_between_states}
\cref{fig:Al_eigen} and \cref{fig:Ti_eigen} display the eigenvectors obtained for Al-substrated and Ti-substrated systems with different compositions at \SI{1033.5}{\kelvin}. To understand the implication of the eigenvectors, informally, one can think of the dynamics of the system as a superposition of many ``principal modes'' with varying timescales.\cite{prinz2011probing}

\begin{figure}[htbp]
    \def\subplotheight{0.2\textwidth}
    \def\verticalskip{0.5\baselineskip}
    \newcommand{\subplot}[2]{
        \sidesubfloat[]{
            \includegraphics[height=\subplotheight]{#1}\label{#2}
        }
    }
    \subplot{figs/eigen_vectors_and_timescales/Al-Al100Ti0-Al-7-1033_dot_5K.png}{fig:Al_eigen_0}
    \subplot{figs/eigen_vectors_and_timescales/Al-Al95Ti05-Al-9-1033_dot_5K.png}{fig:Al_eigen_2}
    \subplot{figs/eigen_vectors_and_timescales/Al-Al90Ti10-Al-7-1033_dot_5K.png}{fig:Al_eigen_3}
    
    \caption{The eigenvectors and associated timescales of Al-substrated dynamic systems with (a) Al$100\%$Ti$0\%$, (b) Al$95\%$Ti$5\%$ and (c) Al$90\%$Ti$10\%$ at 1033.5K.}
    \label{fig:Al_eigen}
\end{figure}

\begin{figure}[htbp]
    \def\subplotheight{0.2\textwidth}
    \def\verticalskip{0.5\baselineskip}
    \newcommand{\subplot}[2]{
        \sidesubfloat[]{
            \includegraphics[height=\subplotheight]{#1}\label{#2}
        }
    }
    \subplot{figs/eigen_vectors_and_timescales/Ti-Al100Ti0-Ti-1-1033_dot_5K.png}{fig:Ti_eigen_0}
    \subplot{figs/eigen_vectors_and_timescales/Ti-Al95Ti05-Ti-9-1033_dot_5K.png}{fig:Ti_eigen_2}
    \subplot{figs/eigen_vectors_and_timescales/Ti-Al90Ti10-Ti-7-1033_dot_5K.png}{fig:Ti_eigen_3}
    
    \caption{The eigenvectors and associated timescales of Ti-substrated dynamic systems with (a) Al$100\%$Ti$0\%$, (b) Al$95\%$Ti$5\%$ and (c) Al$90\%$Ti$10\%$ at 1033.5K.}
    \label{fig:Ti_eigen}
\end{figure}

The eigenvalues $\lambda_i$ determine how slow a process is. In our setup, there are three states, so the number of eigenvalues is $3$. So, in our case, there is with $1 \equiv \lambda_0 > \lambda_1 > \lambda_2$ where $\lambda_0 = 1$ corresponds to the static process. Therefore, for the purpose of analyzing the dynamic processes, we focus on $\mathbf{v}_1$ and $\mathbf{v}_2$. 

Both \cref{fig:Al_eigen} and \cref{fig:Ti_eigen}, for either Al- or Ti-substrated systems with all three compositions of the mixture, exhibit three general patterns of dynamics, as follows:
\begin{itemize}
    \item The values of $\mathbf{v}_1$ imply that the slowest process is the transition between state 1 and 3, with a small portion of the transition between state 2 and 3.
    \item The values of $\mathbf{v}_2$ imply that the second slowest process is the transition between state 2 and 3, with a small portion of the transition between state 1 and 2.
    \item Both identified processes' timescales increase dramatically as the Ti concentration in the mixture increases.
\end{itemize}
The only ``exception'' to this observation is $\mathbf{v}_1$ in \cref{fig:Ti_eigen} (a), which can be easily explained by the fact that the MSM is unable to differentiate between state 2 and 3 in this case, which is already hinted in \cref{fig:Ti-based}. Hence, the $\mathbf{v}_1$ value in \cref{fig:Ti_eigen} (a) does not contradict with the observation.

% eigen_vectors_and_timescales and timescales of the Al and Ti-substrated systems at a lower temperature (\SI{983.5}{\kelvin}) are shown in \cref{fig:Al_eigen_983_dot_5K} and \cref{fig:Ti_eigen_983_dot_5K}. The composition of Al$90\%$Ti$10\%$ is not included in neither Al or Ti-substrated systems at this temperature because the dynamics of atoms at this temperature is so slow and no significant dynamic behavior is observed. From the figures, it is clear that the three patterns of the dynamics observed at \SI{1033.5}{\kelvin} still exists at \SI{983.5}{\kelvin}. Meanwhile, the effect of temperature decrease is significant. By comparison, generally, the identified dynamic processes become roughly two to four times slower, even with only \SI{50}{\kelvin} temperature decrease. For example, for Ti-substrated system with Al$100\%$Ti$0\%$ mixture, the timescale associated with $v_1$ is \SI{5.832}{\nano\second} at \SI{1033.5}{\kelvin}, while the corresponding timescale is \SI{13.68}{\nano\second} at \SI{983.5}{\kelvin}. However, we remark that such comparison is not rigorous as the partitions of the dynamic states are not exactly the same even for the same system at different temperatures. That said, such analysis still provides a new perspective of understanding the properties of the dynamic systems.

Eigenvectors and timescales for the Al and Ti-substrated systems at \SI{983.5}{\kelvin} are included in \cref{asec:timescale_at_983.5K}. It is evident from the results that the three observed dynamic patterns at \SI{1033.5}{\kelvin} persist at \SI{983.5}{\kelvin}. Furthermore, the decrease in temperature has a significant effect as the identified dynamic processes are generally two to four times slower, even with only a \SI{50}{\kelvin} decrease in temperature. For example, in the Ti-substrated system with an Al$100\%$Ti$0\%$ mixture, the timescale associated with $\mathbf{v}_1$ is \SI{5.832}{\nano\second} at \SI{1033.5}{\kelvin} and \SI{13.68}{\nano\second} at \SI{983.5}{\kelvin}. However, we note that such comparisons are not rigorous as the partitions of the dynamic states are not exactly the same even for the same system at different temperatures. Nevertheless, this analysis provides a unique perspective on understanding the properties of such dynamic systems.

% \begin{figure}[htbp]
%     \def\subplotheight{0.2\textwidth}
%     \def\verticalskip{0.5\baselineskip}
%     \newcommand{\subplot}[2]{
%         \sidesubfloat[]{
%             \includegraphics[height=\subplotheight]{#1}\label{#2}
%         }
%     }
%     \subplot{figs/eigen_vectors_and_timescales/Al-Al100Ti0-Al-9-983_dot_5K.png}{fig:Al_eigen_0_983_dot_5K}
%     \subplot{figs/eigen_vectors_and_timescales/Al-Al95Ti05-Al-6-983_dot_5K.png}{fig:Al_eigen_1_983_dot_5K}
    
%     \caption{The eigenvectors of Al-substrated dynamic systems with (a) Al$100\%$Ti$0\%$, (b) Al$95\%$Ti$5\%$ and (c) Al$90\%$Ti$10\%$ at \SI{983.5}{\kelvin}.}
%     \label{fig:Al_eigen_983_dot_5K}
% \end{figure}

% \begin{figure}[htbp]
%     \def\subplotheight{0.2\textwidth}
%     \def\verticalskip{0.5\baselineskip}
%     \newcommand{\subplot}[2]{
%         \sidesubfloat[]{
%             \includegraphics[height=\subplotheight]{#1}\label{#2}
%         }
%     }
%     \subplot{figs/eigen_vectors_and_timescales/Ti-Al100Ti0-Ti-4-983_dot_5K.png}{fig:Ti_eigen_0_983_dot_5K}
%     \subplot{figs/eigen_vectors_and_timescales/Ti-Al95Ti05-Ti-6-983_dot_5K.png}{fig:Ti_eigen_1_983_dot_5K}
    
%     \caption{The eigenvectors of Ti-substrated dynamic systems with (a) Al$100\%$Ti$0\%$, (b) Al$95\%$Ti$5\%$ and (c) Al$90\%$Ti$10\%$ at \SI{983.5}{\kelvin}.}
%     \label{fig:Ti_eigen_983_dot_5K}
% \end{figure}

% \tj{discuss why we do not use Ti-Ti/Al mixture-Al system? 1. simpler system easier to analyze 2. better sampling 3. enforced symmetry (so the actual distribution of states could be more complicated)}
% These results also suggests that we do not need a system with a Ti substrate and an Al substrate. This is because, no matter what substrates are in a system, there is always the liquid Al state, so it can be seen as a reference state which associates the processes. For example, for Al$95\%$Ti$5\%$ mixtures at \SI{1033.5}{\kelvin}, the timescale of the transition between solid Al near Al substrate and liquid Al is \SI{0.553}{\nano\second} where as the timescale of the transition between solid Al near Ti substrate and liquid Al is \SI{8.676}{\nano\second}. From this, one can reasonably infer that in a real world system in which both of the substrates are present, the transition between solid Al near Al substrate and liquid Al is much more rapid than the transition between solid Al near Ti substrate and liquid Al. Also, using substrates of the same composition symmetrize the system, making the system simpler and analysis easier. The results show that for systems with substrates of the same composition generally have three states, which implies that systems with substrates of different compositions will have at least five states (two near-substrate states depending on the type of the substrate and one liquid state in the middle). However, there are at most $7500$ Al atoms on which we sample the dynamics. So, it is reasonable to use substrates of the same composition which results in less states to ensure better sampling of the states.
These findings suggest that we do not necessarily need a system having both Ti and Al substrates to study the dynamic states. This conclusion is drawn from the observation that regardless of the substrates used, the liquid Aluminum state always exists, and can therefore be used as a consistent reference point for all processes. Consider, for instance, the Al$95\%$Ti$5\%$ mixture at 1033.5 Kelvin. The transition from solid to liquid Aluminum near an Aluminum substrate occurs on a timescale of \SI{0.553}{\nano\second}. In contrast, the same phase transition near a Titanium substrate unfolds over a notably longer timescale, approximately \SI{8.676}{\nano\second}. This indicates that in a real-world system where both substrates are present, solid Aluminum transitions to a liquid state much faster near an Aluminum substrate compared to a Titanium one. Furthermore, using substrates of identical composition simplifies the system and makes the analysis easier by creating symmetry. The results also indicate that systems with identical substrates generally have three distinct states. This implies that systems with different substrates will have at least five states -- two near the substrate states that depend on the substrate type and one liquid state in the middle. However, we need to consider that our sample size is limited to a maximum of $7500$ Aluminum atoms. Therefore, it is reasonable to use substrates of the same composition, which results in fewer states, to ensure a more robust sampling of the states.

\subsection{Implications on the diffusive dynamics}
The MD simulation results combined with the MSM analysis shed light on the mechanism of Al diffusion in the Ti--Al system. In a system with pure Ti and Al phases, with the presence of the TiAl\textsubscript{3} grain boundaries, the diffusion of Al at high temperature can be described as three stages. In the initial stage, due to the premelting effect, the Al atoms at the interfaces become liquid, and diffuse into the grain boundaries. This process widens the grain boundaries, opening up channels for Al atoms to move towards the Ti phase, which in turn could enhance the diffusion of Al atoms. In this stage, there is a net flux of Al towards Ti, whereas a very minimal amount of Ti atoms diffuse towards the Al phase. Consequently, there are Al atoms accumulating at the Ti/TiAl\textsubscript{3} interface, continuously opening the grain boundaries along the Ti surface. These are evidenced by the simulation results in \ref{sec:md_results}. Following the initial stage, with the abundance of Al atoms which has opened up the grain boundaries, the Ti atoms will also diffuse into the Al liquid, forming a liquid mixture between the Ti surface and the Al surface. In such a mixture, the transition between the solid Al near the Ti surface and the liquid Al is always slower than the transition between the solid Al near the Al surface and the liquid Al. This essentially means Al atoms in the liquid phase tend to be stably deposited onto the Ti surface more than the Al surface, enhancing the flow of Al atoms going from Al to Ti. Eventually, as the content of Ti increases in the mixture, all the dynamic processes of Al in the mixture become slower, and eventually, the liquid diffusion ceases. Results in \cref{sec:transitions_between_states} show that the timescale of the processes dramatically increase even with just a slight increase in Ti concentration in the mixture. Therefore, the diffusion processes probably will cease when the Ti content in the mixture hit some limit.

% \tj{how does the deposition continue? only by continuously opening the grain boundaries along the Ti surface...}

% Interestingly, the results from \cref{sec:distribution_of_states} and \cref{sec:transitions_between_states} imply a common pattern of the dynamics of Al atoms for both Ti and Al substrate with varying Ti concentration in the mixture. The dynamics is partitioned by three states, which corresponds to increasingly longer distance from the substrates. Regardless of the substrate or the concentration of Ti in the mixture, the slowest process is mostly the transition between state 1 (solid Al) and state 3 (liquid Al), while the second slowest process is mostly the transition between state 2 (interfacial Al) and state 3 (liquid Al). This is intuitive, since the solid phase is the most stable state while the liquid phase is the most active. Remarkably, the MSM provides the timescale of these transitions as a measure of how slow these processes are. Results in \cref{sec:transitions_between_states} show that the timescale of the processes dramatically increase even with just a slight increase in Ti concentration in the mixture. Also, Ti-substrated systems have much larger timescale compared to Al-substrated systems. \tj{keep working on this...}

% The distribution of states from \cref{sec:distribution_of_states} is intuitive. For all systems, it is clear that the dynamics of Al atoms is partitioned by three states, which corresponds to increasingly longer distance from the substrates. This indicates how the dynamics of Al atoms is impacted by the substrate material. 

% \begin{figure}[htbp]
%     \sidesubfloat[]{
%         \includegraphics[width=0.4\textwidth]{figs/Ti-Al100Ti0-Ti.png}\label{fig:Ti-sub0}
%         } \\
%         [\baselineskip]
%     \sidesubfloat[]{
%         \includegraphics[width=0.4\textwidth]{figs/Ti-Al95Ti05-Ti.png}\label{fig:Ti-sub1}
%         } \\
%         [\baselineskip]
%     \sidesubfloat[]{
%         \includegraphics[width=0.4\textwidth]{figs/Ti-Al90Ti10-Ti.png}\label{fig:Ti-sub2}
%         }
%     \caption{The distribution of states of Ti atoms in Ti-substrated system at 1033.5K. The substrates are omitted. (a): Al$100\%$Ti$0\%$; (b): Al$95\%$Ti$5\%$; (c): Al$90\%$Ti$10\%$. The upper panels use a colormap to represents the probability of an atom being in the state. The lower panels plot the corresponding probabilities per atom as a function of the z-coordinate.}
%     \label{fig:Ti-based}
% \end{figure}